\chapter{CONCLUSION AND FUTURE WORK}
\label{ch:conclusion}

\section{Conclusion}

This thesis addressed the automated quantification of bead defects in SEM
micrographs of electrospun nanofibre mats, a measurement currently performed by
hand and therefore slow and difficult to reproduce.

The dataset work is reported in Chapter~\ref{ch:methodology}. Eleven micrographs
of four specimens at three magnifications, carrying 606 manually delineated
beads, were characterised and placed on a physical scale from each micrograph's
burned-in scale bar, and the calibration was validated against the data
themselves: the median equivalent bead diameter agrees to within 15\% across the
three magnifications, which holds only if calibration and annotation are both
consistent.

Two properties of the dataset determine how it must be used. The 143
image--annotation pairs supplied are 11 originals and 132 transformed copies, so
the effective sample size is 11 and the copies cannot cross a partition boundary;
and the eleven micrographs derive from only four specimens, so the specimen
rather than the image is the unit of partitioning. Both are reflected in the
leave-one-specimen-out protocol adopted throughout.

Under that protocol, Mask R-CNN outperforms the mask-predicting alternatives on
every metric measured, reaching $\mathrm{AP}_{50} = 0.298 \pm 0.059$, an
Aggregated Jaccard Index of $0.325 \pm 0.053$ and a mean absolute counting error
of 41.1\%. It reaches neither success criterion of Section~\ref{sec:objectives}:
far short of the $\mathrm{AP}_{50} \geq 0.70$ recorded at the proposal stage, and
two thousandths short of the 0.30 a standard detector reaches on a published
benchmark matched to this object size -- a figure measured on boxes, which the
two detection-only methods here do clear, at 0.305 and 0.315, and which is kept
at the stricter mask level because the objective was stated there. The pipeline
is therefore not yet a replacement for a human operator. Its decomposed Panoptic
Quality -- Segmentation Quality 0.719 against Recognition Quality 0.486 --
locates the deficit: the model delineates well the beads it finds and finds fewer
than half of them, matching 271 of 606 at an overlap of 0.5. Scored on the
micrographs it trained on it reaches 0.437, a gap of 0.139 but still far short of
the criterion, so the shortfall is capacity and training budget rather than
generalisation.

The comparison against the five alternatives answers the question the thesis was
built around, and each of them fails in a way that isolates a different component.

The U-Net semantic baseline merges 48.7\% of ground-truth beads and never splits
one, so it under-counts by 44\% with a bias that never changes sign, while its
mask accuracy differs from Mask R-CNN's by only 0.020 in AJI. What the instance
formulation contributes is therefore not boundary accuracy but the separation of
objects that touch -- and 8.8\% of annotated bead pixels lie in a region claimed
by more than one bead, so touching is common rather than exceptional here.

Faster R-CNN, which is Mask R-CNN with the mask branch removed and everything
else held fixed, counts as well as the full model: 37.8\% against 41.1\%, far
inside the fold-to-fold spread of either, for 15\% less training time. The mask
branch is not what makes the count work; it is what makes the size distribution
measurable at all, and since size in micrometres is one of the two quantities the
laboratory asked for, it is not optional despite contributing nothing to the
count.

YOLOv8 finds more of the real beads than any other method tested -- a recall of
0.696 against Mask R-CNN's 0.447 -- and reports 1223 objects against 606
annotated, for a counting error of 93.2\%. A method can be the best detector
measured and the worst instrument measured, and this one is both.

YOLOv5u, the second model of that family, gives the clearest result of the
comparison and the least usable one. It reaches the highest box
$\mathrm{AP}_{50}$ in the study, $0.315 \pm 0.049$, and the highest
$\mathrm{AP}_S$ of any method, 0.132, and counts three of the four held-out
specimens to within 8.2\%, 15.0\% and 18.3\%, each better than any other method
achieves on any fold. On the fourth it reports 591 objects for 238 beads, an
error of 212.5\%. Re-scoring that fold at the threshold the others selected cuts
the error by half, and no threshold brings it below 34.7\%, because the
specimen's three magnifications need different cuts. The obstacle is not that the
model cannot find these particles but that its confidence scores are not
comparable across magnifications within a specimen -- a solvable problem, and a
different one from the rest of this work.

The classical Otsu-plus-watershed pipeline fails outright at
$\mathrm{AP}_{50} = 0.001$, for a reason measurable in the micrographs: beads are
separated from the surrounding mat by only about 20 grey levels.

Three findings were not anticipated. Detection improves monotonically with
magnification, from $\mathrm{AP}_{50} = 0.227$ at $500\times$ to 0.396 at
$3000\times$, even though the $3000\times$ group holds only 52 of the 606 beads:
object size in pixels, not the number of training examples, is the binding
constraint. The same constraint can be relieved without returning to the
microscope: magnifying the dataset threefold with bicubic interpolation, image and
annotation together, raises $\mathrm{AP}_{50}$ to 0.373 and recall from 0.447 to
0.655, for two and a half times the training cost and no gain in counting. Since
resampling adds no information, the gain is itself the evidence for the diagnosis.
And the naive partitioning protocol did not measurably overstate performance: a
random split over the 11 micrographs gave $\mathrm{AP}_{50} = 0.303$ against 0.298
for the specimen-wise protocol, a difference of 0.005 against a between-fold
standard deviation of 0.059, while AJI and counting error both ran the other way.
It did, however, produce a fold-to-fold spread 1.7 times as wide, which is a
reason to prefer the specimen-wise protocol as the more stable estimator even
though no bias was detected.

A fourth result is negative in a narrower sense. None of the three preprocessing
transforms improved detection: Mask R-CNN scores $\mathrm{AP}_{50} = 0.298$ on
raw micrographs against 0.289, 0.263 and 0.262 with a median filter, background
subtraction and CLAHE, and the contrast measured before any training predicted
the harmful transform correctly and the helpful ones wrongly. One did improve the
laboratory endpoint: background subtraction cut counting error from 41.1\% to
31.3\% by suppressing the runaway over-counts on sparsely populated frames --
the clearest evidence here that detection accuracy and counting accuracy are not
the same objective.

\section{Research Contributions}

\begin{enumerate}
    \item A characterised and physically calibrated dataset of 606 bead defects
    across 11 SEM micrographs of four electrospun specimens, with per-magnification
    pixel sizes recovered from the instrument's own scale bars and validated for
    internal consistency to within 15\%.

    \item An evaluation protocol matched to the hierarchical structure of small
    microscopy datasets, in which the specimen is the unit of partitioning,
    together with a measurement of what a naive split actually costs on such a
    dataset. That measurement returned no resolvable difference in the mean and a
    fold-to-fold spread 1.7 times as wide, which is a more useful result to report
    than the inflation that was expected.

    \item A controlled comparison of six methods -- two-stage instance
    segmentation, the same detector without its mask branch, two one-stage
    alternatives from the same family, semantic segmentation with connected
    components, and classical thresholding -- on identical micrographs under an
    identical protocol. It
    isolates object separation, rather than mask accuracy, as what instance-level
    modelling contributes on this modality, and quantifies it as a 23.8
    percentage point difference in merge rate; and it isolates the mask branch as
    worth every mask-level measurement and nothing at all on the count.

    \item A measurement of preprocessing as an experimental factor rather than as
    a default, on a material whose background carries structure at the same
    spatial scale as the objects, showing that the transform which most improves
    the raw particle-to-mat separation is the one that most harms detection.

    \item A demonstration that the binding constraint on this modality is pixels
    per object rather than the number of annotated objects, obtained by magnifying
    the dataset threefold with bicubic interpolation and holding everything else
    fixed, and separately by scoring each fold's model on its own training
    partition and finding a model that under-fits rather than one that memorises.

    \item An open implementation covering annotation conversion, scale
    calibration, training, evaluation, and the reporting of bead count, areal
        density and size distribution in physical units, in which every table in
    Chapter~\ref{ch:results} is generated from the stored metrics rather than
    transcribed, and whose complete set of results regenerates in 10.8 hours on a
    single freely available GPU.
\end{enumerate}

\section{Limitations}

The limitations are set out in Section~\ref{sec:limitations} and summarised here:
four specimens of a single material system; only 52 annotated beads at the
highest magnification; a single annotator, so annotator bias cannot be separated
from ground truth; a single class, so the distinction between a bead and a
locally thickened fibre segment rests on the annotator's implicit judgement; two
one-stage methods that were not given an input budget matched to the
region-based models, so their accuracy is not a like-for-like architectural
comparison; and a single training run per configuration at a fixed seed, so every
figure is a point estimate and the reported spreads describe sensitivity to the
held-out specimen rather than run-to-run variability.

\section{Future Work}
\label{sec:futurework}

\subsection{Short-term}

\textbf{Address the recognition deficit directly.} A Segmentation Quality of
0.719 against a Recognition Quality of 0.486 says effort should go into finding
beads rather than refining masks. The first steps are a longer schedule than the
1500 iterations the available GPU time allowed, a hyperparameter search this work
did not perform, and a loss that up-weights recall; the train-versus-test
comparison of Section~\ref{sec:generalisation} argues this is the right place to
spend, since the model scores 0.437 on data it has already seen. YOLOv8's recall
of 0.696 gives the target, and the magnified run of
Section~\ref{sec:upsample_results} reaches 0.655 with the same architecture as
the baseline's 0.447.

\textbf{Calibrate the score threshold per magnification, not per fold.} This is
the cheapest experiment suggested by any result here and the one with the largest
measured headroom. YOLOv5u's single catastrophic fold is mostly an operating
point: one threshold was selected for a specimen whose three magnifications need
different ones, and that fold's error halves at the value the other folds chose.
Selecting per magnification group, still on training data only, is a few lines of
change to the existing selection code and costs no additional training.

\textbf{Suppress false positives rather than chase recall.} The complementary
reading of the same comparison is that YOLOv8's recall advantage is bought with
801 spurious detections. A post-hoc classifier over detections would test whether
the recall of one method and the precision of another can be had together.

\textbf{Extend the specimen count.} The most valuable addition to the data is
more specimens rather than more images of the same ones. Four support four folds;
ten would support a materially stronger claim, at an annotation cost per specimen
this work has already measured, and would give the protocol comparison of
Section~\ref{sec:leakage} the resolution these four folds lacked.

\textbf{Measure inter-annotator agreement.} A second annotator re-delineating a
subset would convert the largest stated limitation into a measured quantity and
establish the ceiling against which any model should be judged. At
$\mathrm{AP}_{50} = 0.298$ it is not known how much of the shortfall is model
error and how much is disagreement a second human would also show, and a model
that matches inter-annotator agreement has reached the useful limit of the task.

\textbf{Use the unannotated micrographs.} If further micrographs exist without
annotations, they can be used for self-training or consistency-based
semi-supervised learning, both of which are well suited to a regime with abundant
images and scarce labels.

\textbf{Shape-prior architectures.} Bead defects are approximately ellipsoidal and
convex, which is exactly the assumption behind StarDist
\cite{schmidt2018stardist} and the flow-based representation of Cellpose
\cite{stringer2021cellpose}. Both are designed for the touching-object problem
that Section~\ref{ch:results} identifies as the central difficulty, and are a
natural extension of the comparison.

\subsection{Long-term}

\textbf{Image at the magnification the measurement needs.} The magnification
result suggests a change in acquisition practice rather than in modelling. Higher
magnification gives more reliable per-bead measurements while a count is only as
good as the number of objects in the field, so a mixed protocol -- a wide field
for density, several high-power fields for size -- may measure both better than
any single setting. Testing it needs micrographs acquired for the purpose, which
would also give the counting comparison across magnifications the sample size
three or four micrographs per group denied it here.

\textbf{Joint fibre and bead analysis.} Bead quantification is one half of the
morphological characterisation of a nanofibre mat; fibre diameter distribution is the
other, and tools such as DiameterJ \cite{hotaling2015diameterj} address it
separately. A single model reporting both from one micrograph would cover the
characterisation a laboratory actually needs.

\textbf{From measurement to process feedback.} If bead statistics could be linked to
the spinning parameters that produced them, the pipeline would move from
characterisation towards prediction -- estimating, from a micrograph, which parameter
is out of range. This requires micrographs labelled with their process conditions,
which the present dataset does not carry.

\textbf{Annotation assistance from foundation models.} The Segment Anything Model
\cite{kirillov2023segment} cannot recognise beads on its own, but prompted with a
detector's boxes it could refine mask boundaries, or accelerate the annotation of
new specimens by turning a click into a polygon. Given that annotation effort is the
binding constraint on this problem, this is the intervention with the highest
leverage.

\section{Final Remarks}

The pipeline built here does not yet replace the operator with the mouse. At an
$\mathrm{AP}_{50}$ of 0.298 and a per-micrograph counting error that ranges from
$-75\%$ to $+119\%$, no laboratory should report its output on a single field as a
measurement without checking it.

What the work does establish is where the difficulty lies, which is worth more at
this stage than a higher number would have been. It is not segmentation accuracy:
the masks overlap their targets by 0.72 on average. It is not the size
distribution: the predicted median diameter agrees with the manual measurement to
3.2\%. It is detection of small, low-contrast objects that touch one another, and
the measurements that isolate it are the recognition-versus-segmentation split of
Panoptic Quality, the merge and split rates that separate the instance model from
the semantic one, the magnification trend, and the threefold resampling that
reproduces that trend without adding a single new observation.

One methodological result deserves to outlive the specific numbers. Across this
comparison, the ranking of the methods depends on which quantity is asked for.
YOLOv8 beats the U-Net on average precision and loses to it on counting. YOLOv5u
posts the best detection scores in the study and a counting error that is either
the best measured or the second worst, depending on which specimen is held out.
The two methods with the best agreement to the manual size distribution, by the
Kolmogorov--Smirnov statistic, are the two that over-count most severely. A study
of this problem that reported a single accuracy figure would have ranked the
methods wrongly for the laboratory that has to use them, and would not have been
able to tell.

Three of the results reported here are negative: the accuracy criterion was not
met, the data-leakage effect the protocol was designed to prevent did not appear,
and no preprocessing transform improved detection. All are reported as measured
rather than reframed, and the second in particular is the more useful finding,
because it tells the next person working on a dataset of this shape that the
expensive precaution matters for the variance of their estimate rather than for
its mean. A thesis that reported only the comfortable half of its results would
have given a worse account of what this dataset can and cannot support.
